\documentclass[10pt,a4paper]{article}
\usepackage[margin=18mm]{geometry}
\usepackage[T1]{fontenc}
\usepackage{lmodern,graphicx,booktabs,amsmath,amssymb,array}
\usepackage[hidelinks]{hyperref}
\setlength{\parindent}{0pt}
\setlength{\parskip}{5pt}
\begin{document}
{\Large\bfseries The actual degree-four cover has four distinct branch values}\par
{\large\bfseries A canonical ramification line and exact quadratic-field images}\par
24 November 2026 \quad crypto PR5110 \quad [pcc:unverified]

\textbf{Result and scope.} The actual global map $f_4:C\to E$ now has
a complete ramification divisor and every branch image, including the
quadratic orbit. Exact native coefficients, local differential values,
factor identities and extension-field evaluations pass strict checks.
The interpretation uses the named actual-map, smoothness, canonical-basis
and tame Riemann--Hurwitz bridges. Those remain written mathematics,
with full geometric formalization and independent signed standing open.
Native B 2bf624, native C 147cd6; written proofs 604246 and de29b0.

\[
\begin{split}
C/\mathbb F_7:\;&X^4+4XY^3+2Y^4+5X^2Z^2+6XYZ^2+2Y^2Z^2+Z^4=0,\\
E/\mathbb F_7:\;&v^2=u^3+2u+6.
\end{split}
\]

\textbf{The actual differential fixes the line.} On $X=1$, write
$y=Y/X,z=Z/X$ and $\theta=dy/F_z=-dz/F_y$. Exact chart transitions
and genus three give the genuine basis $\theta,y\theta,z\theta$ of
holomorphic differentials. All 17 inherited actual local values give
the complete rank-three linear certificate
\[
 f_4^*\!\left(\frac{du}{2v}\right)=(2+z)\theta,\qquad h=2X+Z.
\]
Existence comes from the actual holomorphic pullback; equality at 17
points determines its linear form, since a nonzero line meets the
quartic in degree four. Rational zero points alone are not used to
guess the complete ramification divisor.

\textbf{The entire line section.} Parametrize $h=0$ by $[U:V:5U]$.
Full homogeneous multiplication yields
\[
\begin{aligned}
F(U,V,5U)&=2U^4+3U^3V+U^2V^2+4UV^3+2V^4\\
 &=2(V-4U)(V-5U)(V^2+4UV+6U^2).
\end{aligned}
\]
The $V^4$ coefficient 2 excludes an infinity zero. All seven nonzero
evaluations prove the quadratic irreducible; the factors are distinct
and separable. All four geometric zeros are simple. The true degree-four
map is separable and tame in characteristic seven, so their indices
are $e=2$, and their total ramification degree is four.

\begin{center}
\includegraphics[width=\linewidth,height=83mm,keepaspectratio]{branches.pdf}
\end{center}
Here $\alpha^2+4\alpha+6=0$ and $\alpha^7=3-\alpha$.
The full native equality matrix of the four target values is the
identity: there are four distinct branch values, with the displayed
Frobenius action.

\newpage
{\Large\bfseries Actual quotient-field evaluation and reproducible evidence}\par

\textbf{The quadratic orbit is evaluated from the functions.} Use the
observed field $K=\mathbb F_7[\alpha]/(\alpha^2+4\alpha+6)$, with the
exact basis $1,\alpha$. For the actual parent $g=f_4-(3,5)$ and its
homogeneous functions $g=(p_4/q^2,p_6/q^3)$, at $(1:\alpha:5)$ the
native certificate retains
\[
\begin{gathered}
q=3+5\alpha,\qquad p_4=2+\alpha,\qquad p_6=1,\\
g=(5-\alpha,3-2\alpha),\qquad f_4=(3-\alpha,4-\alpha).
\end{gathered}
The second point and image are the genuine Frobenius conjugates.
Both curve and elliptic equations hold; every actual denominator is
nonzero. All 49 field Frobenius/norm records, conjugation-squared and
Frobenius-49 identities, and 48 norm inverses are retained. These are
operations in the actual irreducible polynomial quotient, not integers
modulo 49. Conrad's finite-fields Theorems 1.1, 2.5, 4.1 and 4.6 provide
the exact field and conjugation hypotheses.

\textbf{The base-field branch description.} The complete support is
$(3,2),(3,5)$ and the degree-two closed point on the affine target chart
\[
u^2+4u+6=0,\qquad v=u+1.
\]
The monic quadratic and linear relation are derived from the computed
coordinates, then checked on both values. The residual pair's elliptic
\emph{group} sum is $(5,1)$, separately computed and Frobenius-fixed.
Since the rational pair sums to $O$, the full branch divisor class
$B-4O$ is represented by $(5,1)$ under the standard elliptic Abel/group
identification. This written bridge is distinct from coordinate field trace.

\begin{center}
\begin{tabular}{lllll}
\toprule
New unit & Frozen source & Whole lock files & Strict phases & Full stdout\\
\midrule
Canonical divisor B & \texttt{3c6d3b26} & 51 & $7/7$ pass & 3915 bytes\\
Residual branches C & \texttt{c17417cd} & 72 & $7/7$ pass & 6467 bytes\\
\bottomrule
\end{tabular}
\end{center}
Each unit used complete GitHub readback before first calls, pinned
Rust 1.99.1, warnings errors, opt3, overflow/debug assertions,
120 seconds/2 GiB/one worker, and unchanged source/tool/qualifier guards.
Each retains 40 whole small originals. B keeps every linear row,
RREF/residual and full factor witness. C derives both residual points
from the actual observed factor; no distinct-image pattern was assumed.
The old 17-point image population was not repeated.

\textbf{Real controls and first failure.} B's actual linear and binary
coefficient mutations reject. At each residual point C changes the actual
$p_4$ $X^4$ coefficient by one; both changed images fail the elliptic
equation and actual-map equality. Pristine restores. First A compiled
successfully but strict Clippy rejected an indexed two-column loop;
no mathematical case or qualifier execution ran. Its 33 whole originals,
666-byte diagnostic and compiled-only ELF remain unchanged. B fixes
iterator spelling with identical mathematical predicates and no suppression.

\textbf{Original custody and CI.} Original B/C executed and separately
copied0555 ELFs are 4648624/4638624 bytes, hash prefixes
\texttt{549fcb54}/\texttt{7dd3c67e}; A's compiled-only original is
4648768 bytes, \texttt{2dffbb14}. Whole originals remain in RAM pending
approved durable custody. The focused native-only CI setup passes syntax
and actionlint; new CI artifacts remain distinct. Required/applicable CI,
stacked parents and the frozen main ledger fork retain their merge rules.

\textbf{Next concrete chain.} Use the actual branch divisor and class to
review and construct $S_4$ sign/cubic resolvents, then test whether their
natural Jacobian correspondences carry a component beyond the common
elliptic factor. No such nontrivial Prym transfer is assumed. Full sources,
locks, proof maps and inventories are in the study. Primary geometry:
Vakil21.5.B; Stacks0C1B and0BYA. All scientific standing remains
[pcc:unverified].
\end{document}
